\chapter{Discussion}
\label{ch:discussion}

\section{Interpretation of Results}
\label{sec:interpretation}

\subsection{Entity Classification Success}

The LayoutLMv3 V4 entity classifier achieves F1~$= 0.827$, demonstrating that
the layout-aware encoder effectively distinguishes key, value, and other tokens
using text and layout representations. This result indicates strong entity
classification in the reported setup: by
constraining predictions to existing OCR tokens, the model prevents the
generation of text absent from the OCR input. It can still falsely label an
existing token as a key or value.

\subsection{Token-Level Linking Failure}

The near-zero link F1 ($0.008$) of V1 is the central negative result of this
thesis. The failure is consistent with severe token-level class imbalance and a
mismatch between token-level candidates and span-level relations.
With $\sim$450:1 negative-to-positive ratio at the token level, the BCE loss
is dominated by negative examples even with \texttt{pos\_weight} correction.
This is consistent with SPADE~\cite{spade} and BROS~\cite{bros}, which both
operate at the segment/span level rather than the token level.

\subsection{Comparison to Baselines}

Under the released pair-level macro benchmark, reconstructed Mistral reaches
regular-KVP F1 $=0.598$ for text and $0.521$ for text+location. The published
regular-KVP values are $0.659$ and $0.611$. The initial V2 span-level linker
fails under the internal pooled diagnostic, remaining near V1's zero score.
The diagnostic analysis (Section~\ref{sec:res567}) identified three compounding
data-pipeline bugs that prevented a clean assessment of the span-level linker
in V2. Once corrected, V4 reaches regular-KVP F1 $=0.334$ for text
and $0.309$ for text+location, below both reconstructed Mistral and the
published baseline. Its All score is lower ($0.214/0.199$) because its decoder
emits only regular pairs. Internal diagnostic analyses
(Section~\ref{sec:linker-diagnostics}) indicate that relation modeling is the
principal remaining internal bottleneck rather than entity detection or threshold choice;
a decode-time recovery of unlinked spans (Analysis~C) cheaply extends coverage
to unkeyed and unvalued at lower accuracy, raising All F1 to $0.257/0.237$
without retraining.

\subsection{Per-Cluster Performance}

Stratifying Stage 3 results by the Stage 2 annotation-geometry
clusters (Table~\ref{tab:exp1-cluster}) shows an association between
annotation-box geometry and performance. The generative baseline reaches regular-pair text
F1 of 0.651 in the higher-count, broad-spread cluster and 0.498 in the
lower-count, compact cluster; combined F1 is 0.569 and 0.429, respectively.
A separate pooled error audit gives unsupported-prediction rates of 37.1\%
and 53.0\%. These clusters do not measure OCR density or document genre, so
the difference should not be described as a failure on a semantic document
type. The discriminative LayoutLMv3 model cannot generate text beyond the OCR
input because it only labels tokens that exist on the page.

\begin{table}[ht]
  \centering
  \caption{Stage~4 regular-KVP macro F1 by page-level annotation-geometry cluster.}
  \label{tab:stage4-cluster}
  \begin{tabular}{@{}lrrrr@{}}
    \toprule
    \textbf{Stratum} & \textbf{Assigned} & \textbf{Scored} & \textbf{Text F1} & \textbf{Text+location F1} \\
    \midrule
    Cluster 0 (higher count, broad spread) & 319 & 304 & 0.290 & 0.264 \\
    Cluster 1 (lower count, compact)       & 213 & 179 & 0.406 & 0.383 \\
    Geometry unavailable                    &  49 &   0 & ---   & ---   \\
    \bottomrule
  \end{tabular}
\end{table}

V4 performs better on Cluster~1 under both regular-pair modes, the reverse of
the Stage~3 ordering. This is an association over annotation geometry, not
evidence that a semantic layout class is generally easier. The 49 pages with
unavailable geometry have empty regular ground truth and are excluded by the
official macro evaluator. Overall benchmark scores remain unchanged because
only the test-page grouping was recomputed.

\section{Root-Cause Analysis}
\label{sec:root-cause-analysis}

The progression from V1 to V4 illustrates a common but underreported pattern
in structured prediction pipelines: architectural improvements cannot
compensate for corrupted supervision.  The diagnostics identified three
interacting pipeline failures affecting span construction and relation
supervision.  Because V4 introduced the corrections jointly, their separate
numerical contributions were not measured; the following mechanisms are
described without ranking individual effects.

\paragraph{Missing teacher forcing.}
The linker was trained on spans derived from noisy \texttt{argmax} predictions
rather than ground-truth labels.  Span boundaries were inconsistent with the
ground-truth link labels, which could corrupt span boundaries and create a
train--test mismatch that impeded convergence.

\paragraph{Bounding-box filter.}
The strict $(x_1 < x_2) \wedge (y_1 < y_2)$ predicate rejected tokens whose
normalised coordinates had $x_1 = x_2$ after LayoutLMv3's integer
quantisation to $[0, 1000]$.  This removed many valid tokens---49\% of
predicted key tokens and 46\% of value tokens---halving the number of entity
spans available for linking.

\paragraph{Cross-product link labels.}
The label generator created all-pairs links between key and value words rather
than a single representative link per KVP.  Combined with OCR blocks that
concatenated multiple words into single ``word'' entries, this corrupted the
relation-supervision signal: the pipeline produced $\sim$1{,}000 positive link
entries per document instead of $\sim$8, making the span-level collapse
produce near-uniform positive labels.

\noindent The cross-product label issue was invisible at the token level
(V1), where the dense $n_k \times n_v$ matrix already had many positives.
It became observable only when the span-level collapse was introduced in V2,
because the small $K \times V$ matrix should have been sparse but was not.
The span-level reformulation therefore served as a diagnostic
tool, not only an architectural improvement.

\section{Error Analysis}
\label{sec:error_analysis}

\subsection{Entity Classification Errors}

Entity errors fall into three categories: (i)~boundary errors, where partial
spans are predicted (e.g., missing the last token of a multi-word key);
(ii)~type confusion, where a key is predicted as a value or vice versa; and
(iii)~false positives on non-entity tokens that appear near key-value regions.

\subsection{Linking Errors (V1)}

V1 linking errors are dominated by false negatives: the model predicts only
200 links against 3{,}776 ground-truth links. Of the 200 predictions, 15
are correct (7.5\% precision). The model adopts an overly conservative
strategy, preferring to predict no link rather than risk a false positive ---
a rational response to the extreme class imbalance in the training signal.

\subsection{V2 Linking Errors}

V2 span-level linking achieves link F1\,=\,0.007, marginally below V1's
0.008.  The model predicts 498~links across 581 test documents, of which
37 are correct (7.4\% precision).  Recall is 0.3\%, with only 37 of
3{,}776 ground-truth links recovered.

The error pattern differs from V1 in a revealing way: the linker produces
more predictions (498~vs.\ 200) but no improvement in true positives.  The
additional predictions are false positives caused by the corrupted
link supervision (Bug~3), which trained the linker to associate nearly any
span pair.  With the V4 corrections, the linker receives correct sparse
supervision where each KVP produces exactly one positive span link.

\section{Relation-Modeling Context and Experimental Scope}
\label{sec:relation_context}

The experimental comparisons in this thesis remain anchored to the KVP10k
Mistral baseline and the LayoutLMv3 encoder with a biaffine linker. As discussed
in Section~\ref{sec:experimental_scope}, other relation-extraction methods were
available before the project started, but implementing them would have required
separate model designs and full experimental pipelines. SPADE~\cite{spade},
BROS~\cite{bros}, GeoLayoutLM~\cite{geolayoutlm},
KVPFormer~\cite{kvpformer}, and LayoutPointer~\cite{layoutpointer} are therefore
included as methodological context and future alternatives rather than as
additional experimental baselines.

SPADE is reported on CORD, while BROS is evaluated on several
document-understanding benchmarks, including FUNSD, CORD, SROIE, and SciTSR;
neither work is evaluated on KVP10k. Direct numerical comparison is therefore
not possible.
Their segment/span-level formulations provide context for the V2 design and
its reduction of the candidate-space imbalance observed in V1, but do not make
the reported results directly comparable across benchmarks.

\section{Limitations and Future Work}
\label{sec:limitations}

\subsection{Current Limitations}

Three limitations of the present work should be noted. First, our training set
covers only 55.8\% of the unique KVP10k pages, owing to PDF download failures
and scanned documents lacking a text layer; an OCR-based pipeline (for example
Tesseract) would increase coverage. Second, the reconstructed Mistral baseline
used LoRA rank $r = 4$; the published work does not disclose enough
adapter and training detail to reproduce its configuration exactly, which limits
strict comparability between Stage~3 and the published baseline. Third, the
current pipeline processes a single page at a time, so multi-page documents
with cross-page key--value pairs cannot yet be handled.

\subsection{Future Directions}

Several directions follow naturally from these findings. The most promising is
to replace the linker head with a segment-graph relation decoder such as
SPADE~\cite{spade} or BROS~\cite{bros}: within the current biaffine-linker
pipeline, our oracle analysis (Section~\ref{sec:linker-diagnostics}) identifies
relation modeling as the main remaining bottleneck, making linkers that model
directed key$\to$value dependencies explicitly the natural next step. This was
scoped as future work rather than implemented, since it requires
a new decoder head and full retraining. Relation-focused designs from
GeoLayoutLM~\cite{geolayoutlm}, KVPFormer~\cite{kvpformer}, and
LayoutPointer~\cite{layoutpointer} are additional alternatives to the current
biaffine linker. Further directions include handling
multi-page documents with cross-page linking, integrating end-to-end OCR to
increase data coverage, evaluating cross-lingual performance on XFUND~\cite{xfund}, and
exploring contrastive pre-training for the linker head.
